\section{Executable reference artifacts}
\label{app:reference}

The archive's \code{reference/scope\_reference.py} is a standard-library model
of the finite algorithms.  It includes:

\begin{itemize}
\item an AST and evaluator for equality, order, Boolean connectives, and nested
  bounded quantifiers;
\item least bounded witness search and a minimality certificate;
\item monotone capped demand closure with a runtime rank check;
\item a finite DAG checker for sealed, relational, and named escape modes;
\item exact observer independence and metric diameter checks;
\item a JSON demonstration generator.
\end{itemize}

The bundled unit suite has sixteen tests.  It includes two nested bounded
quantifiers, implication truth tables, parameter environments under binders,
least search with and without a witness, a cyclic demand system, rejection of
an unjustified exact name, a missing observer-coherence proof, and exact and
metric observer examples.

Run from the archive root:

\begin{verbatim}
python3 -m unittest -v reference/test_scope_reference.py
python3 reference/scope_reference.py --pretty > examples/demo.json
\end{verbatim}

Passing these tests shows that the reference algorithms and metadata checks
behave as documented on the examples.  It does not prove the universal iPRA
reflection or compiler theorems.  Those are certified only when the Lean
modules construct accepted \code{Prf PRAAx} terms.
